\scanpage{26}
\noindent\hfill ALGEBRE\quad COMMUTATIVE\hfill Rédaction n\textsuperscript{o}355

CHAPITRE IV - DECOMPOSITION PRIMAIRE

\uline{Décomposition primaire et extension des scalaires} (Etat 2)

\begin{center}-:-:-\end{center}

L'état 1 porte le n\textsuperscript{o}352. Après l'avoir lu (et même avant), Grothendieck a signalé au rédacteur qu'on pouvait généraliser le résultat principal, en remplaçant l'un des anneaux par un module. La démonstration s'étend sans changement, et c'est bien plus joli. Le rédacteur en a profité pour enlever quelques hypothèses noethériennes inutiles.

\begin{center}-:-:-\end{center}

Dans tout ce qui suit, on désigne par $A$ et $B$ deux anneaux commutatifs, et l'on considère un homomorphisme $\handwritten{\pi}:A\to B$. Cet homomorphisme fait de $B$ une $A$-algèbre. On se donne en outre un $B$-module $F$ que l'on suppose \uline{plat sur} $A$.

PROPOSITION 1.- \uline{Soit} $\underline p$ \uline{un idéal premier de} $A$. \uline{Pour tout} $\underline P\handwritten{\in}\operatorname{Ass}_B(F\handwritten{/}\underline p F)$ \uline{on a} $\handwritten{\pi}^{-1}(\underline P)=\underline p$.

Puisque $\handwritten{\pi}(\underline p)$ annule $F\handwritten{/}\underline p F$, on a $\handwritten{\pi}(\underline p)\handwritten{\subset}\underline P$, d'o\handwritten{\`}{u} $\underline p\handwritten{\subset}\handwritten{\pi}^{-1}(\underline P)$. D'autre part, si $a\handwritten{\in} A-\underline p$, l'homothétie de rapport $a$ est injective dans $A\handwritten{/}\underline p$, donc aussi (puisque $F$ est $A$-plat) dans $F\handwritten{/}\underline p F$, ce qui montre que $\handwritten{\pi}(a)\handwritten{\notin}\underline P$, CQFD.

PROPOSITION 2.- \uline{Soit} $E$ \uline{un} $A$-\uline{module}. \uline{On a} :

\[
\operatorname{Ass}_B(E\handwritten{\otimes}_A F)\handwritten{\supset}\handwritten{\bigcup}_{\underline p\handwritten{\in}\operatorname{Ass}_A(E)}\operatorname{Ass}_B(F\handwritten{/}\underline p F)
\tag*{\makebox[0pt][r]{\makebox[\linewidth][l]{$(\ast)$}}}
\]

\uline{Lorsque} $A$ \uline{est noethérien}, \uline{l'inclusion} $(\ast)$ \uline{est une égalité}.

Soit $\underline p\handwritten{\in}\operatorname{Ass}_A(E)$. Par définition de $\operatorname{Ass}_A(E)$, il existe une suite exacte :

\[
0\to A\handwritten{/}\underline p\to E\,.
\]

Puisque $F$ est plat sur $A$, on en déduit une suite exacte :

\[
0\to F\handwritten{/}\underline p F\to E\handwritten{\otimes}_A F\,,
\]

% archive sheet number: 25
